\chapter{Keeping the Log Small}
\label{ch:logkeep}

A write-ahead log only grows. Every commit appends, nothing removes, and a
database that runs for months would keep every record since its creation and
replay from the beginning of time on every crash. ARIES as described in
\cref{ch:aries} is correct but not yet livable. This chapter covers the three
pieces of housekeeping a long-lived page store needs: cutting the log after a
checkpoint, making page records smaller, and checkpointing without stopping
the world. Each is a statement about what recovery can still need, made
precise.

\section{What recovery still needs}

After a crash, ARIES redo starts at the smallest recLSN in the dirty page
table: the first change to any page not yet written. Undo walks back through
every loser transaction, so it needs each active transaction's records back
to its first one. Everything older than both is dead weight:
\[
\mathit{keep} \;=\; \min\Big(\min_{p\ \text{dirty}} \mathit{recLSN}(p),\ \ \min_{T\ \text{active}} \mathit{firstLSN}(T)\Big).
\]
A record below \textit{keep} describes a change already on disk, made by a
transaction already finished; no crash can need it. A checkpoint is the moment
these two minima are known and written down, which is why truncation follows a
checkpoint.

\begin{keyidea}
The log is a queue of obligations: redo for pages not yet written, undo for
transactions not yet finished. Its tail can be cut exactly at the oldest
obligation still open.
\end{keyidea}

\section{Cutting the log}

\Cref{ch:durability} empties its log after a snapshot, because the snapshot
holds everything. The page store cannot empty it: after a checkpoint, dirty
pages and active transactions still depend on records near the end. It keeps
the suffix from \textit{keep} on, and must do so without ever leaving a log
that recovers differently:

\begin{enumerate}
\item Write the records from \textit{keep} on to \code{wal.tmp}, and \code{fsync} it.
\item \code{rename} it over \code{wal}, and \code{fsync} the directory.
\end{enumerate}

\begin{center}
\begin{tabular}{@{}ll@{}}
\toprule
Crash after & On reopen \\
\midrule
writing \code{wal.tmp} & the old, full log: recovery as before \\
the rename & the cut log: same recovery, less to read \\
\bottomrule
\end{tabular}
\end{center}

LSNs are counters, not file offsets, so the records keep their LSNs and every
pageLSN on disk still points where it did. The file only needs to know the
LSN of its first record, which goes in its header.

\begin{decision}{Truncate by copying the tail (planned)}
\dpart{Choice} After a checkpoint, copy the needed suffix to a new file and
rename it over the log.
\dpart{Alternatives} A log made of fixed-size segment files, where truncation
deletes whole old segments, as PostgreSQL and InnoDB do.
\dpart{Why} One file keeps the log manager, the scanner, and the crash harness
as they are, and copying is cheap right after a checkpoint, when the suffix is
short.
\dpart{Cost} Copying a long suffix while a long transaction stays open; segment
files would avoid the copy and cost a directory of files to manage.
\end{decision}

\section{Torn pages and smaller page records}

\Cref{ch:aries} logs a page's full image with every change. That is simple and
robust, and it makes the log as large as the pages it touches: insert one
24-byte key into a 4\,KiB leaf, and the log grows by 4\,KiB.

Logging only the changed bytes is the obvious saving, and it has one hazard.
A page write is not atomic: the device may write some sectors of a page and
not others, leaving a \term{torn page} that is neither the old version nor the
new. Redo that applies a byte diff to a torn page produces garbage, because the
diff assumes the bytes around it are the old ones.

PostgreSQL's answer, \term{full-page writes}, is the one the plan adopts: the
\emph{first} change to a page after a checkpoint logs the page's full image,
and later changes until the next checkpoint log only the byte ranges that
differ from the page as it was when latched. Recovery starts from the
checkpoint, so for every page redo touches, it meets that page's full image
first and rebuilds the page whole, torn or not; every later diff then applies
to a known image.

\begin{center}
\begin{tabular}{@{}lr@{}}
\toprule
Inserting 100 keys into one leaf, one checkpoint & Log bytes \\
\midrule
Full image every change & $100 \times 4{,}096 \approx 410$\,KB \\
Full image first, then diffs (about 40 bytes each) & $4{,}096 + 99 \times 40 \approx 8$\,KB \\
\bottomrule
\end{tabular}
\end{center}

A diff must also say what it applies to. Redo applies a diff only to a page
whose pageLSN equals the LSN of that page's previous record, so a diff can
never land on the wrong version of the page.

\begin{decision}{Full image after a checkpoint, then diffs (planned)}
\dpart{Choice} Log a page's full image on its first change after each
checkpoint, and byte-range diffs for later changes.
\dpart{Alternatives} Full images always (today); diffs always, with a
doublewrite buffer to repair torn pages (InnoDB).
\dpart{Why} Torn pages are repaired from the log with no second write path,
and the log shrinks by roughly the ratio of a page to a change.
\dpart{Cost} Every checkpoint makes the next change to each page expensive
again, so checkpoints that are too frequent waste log; the benchmark measures
log bytes per triple before and after.
\end{decision}

\section{Fuzzy checkpoints}

The sharp checkpoint of \cref{ch:aries} holds every statement out, writes
every dirty page, and then records that nothing is dirty. Its pause grows with
the buffer pool. A \term{fuzzy checkpoint} writes no pages at all:

\begin{enumerate}
\item Append a \textbf{begin-checkpoint} record.
\item Copy the transaction table and the dirty page table, each page with its
recLSN, while transactions keep running, and append them in an
\textbf{end-checkpoint} record.
\item Force the log, and record the begin-checkpoint's LSN in a known place.
\end{enumerate}

Recovery's analysis pass starts from the begin-checkpoint record with the
tables the end record carries, and redo starts from the smallest recLSN in
the dirty page table. The checkpoint's copy of the tables may be slightly out
of date by the time it is written, which is why it is ``fuzzy'': analysis reads
the records between begin and end and brings the tables up to date, just as
it does for every record after the checkpoint.

Writing dirty pages becomes a background job (a page cleaner, in InnoDB's
terms) that keeps the smallest recLSN moving forward, so that redo after a
crash, and the log that truncation must keep, stay short. The two mechanisms
work together: the cleaner advances \textit{keep}, and the checkpoint records
where it is.

\begin{figure}[ht]
\centering
\begin{tikzpicture}[font=\scriptsize,rec/.style={draw,minimum height=0.5cm,inner sep=2pt,fill=white}]
\node[rec] (a) {upd T5 p12};
\node[rec,right=0pt of a,fill=boxbg] (b) {begin ckpt};
\node[rec,right=0pt of b] (c) {upd T6 p40};
\node[rec,right=0pt of c,fill=boxbg] (d) {end ckpt: ATT \{T5, T6\}, DPT \{p12: 88, p40: 97\}};
\node[rec,right=0pt of d] (e) {commit T5};
\draw[decorate,decoration={brace,mirror,amplitude=4pt}] ([yshift=-3pt]b.south west) -- node[below=4pt]{transactions keep running} ([yshift=-3pt]d.south east);
\end{tikzpicture}
\caption{A fuzzy checkpoint. No page is written; the end record carries the
active transactions and each dirty page's recLSN as they were seen.}
\end{figure}

\begin{status}
Planned (plan 11.1--11.3), after the page store reaches the shell
(\cref{ch:pagestore}). Today the log is never truncated, every page record is a
full image, and the checkpoint is sharp. The sharp checkpoint stays as an
option for comparison, and \code{benchmarks/logging/} will measure log bytes per
inserted triple, the checkpoint pause, and recovery time.
\end{status}

\begin{checkyourself}
\item The dirty page table holds p3 (recLSN 140) and p9 (recLSN 212); T7 is active with first LSN 120 and T8 with first LSN 300. From which LSN must the log be kept?
\item Why may truncation keep the old LSNs instead of renumbering from zero?
\item What goes wrong if redo applies a byte diff to a torn page, and how does a full image after each checkpoint prevent it?
\item Why does a diff record need the LSN of the page's previous record?
\item What does a fuzzy checkpoint write, and why may its tables be out of date when they reach the log?
\item Why does a fuzzy checkpoint make the page cleaner more important?
\end{checkyourself}

\section*{Further reading}
Mohan et al.\ \cite{aries} on fuzzy checkpoints and the dirty page table;
Gray and Reuter \cite{grayreuter} on log management; the PostgreSQL
documentation on WAL and full-page writes \cite{pgwal}.
